\documentclass[10pt]{article}
\usepackage[T1]{fontenc}
\usepackage[utf8]{inputenc}
\usepackage[french]{babel}
\usepackage{amsmath,amssymb,mathrsfs,mathpazo}
\usepackage{geometry}
\geometry{a4paper,margin=2cm}
\usepackage{graphicx,array,tabularx,tabularray,colortbl,multirow,multicol,booktabs,enumitem,pifont,ulem,fancybox,diagbox,cancel}
\usepackage{pstricks,pstricks-add,pst-plot,pst-tree,pst-node,pst-text,pst-func,pst-math,pst-bspline,pst-3dplot,pst-eucl,pst-solides3d}
\let\pscalculate\undefined
\usepackage{pst-bezier}
\usepackage{tikz,pgfplots,tkz-tab}
\usetikzlibrary{arrows,arrows.meta,patterns,calc,positioning,trees,decorations.markings,shapes.geometric,shapes.misc,shapes.arrows,intersections,decorations.pathmorphing,decorations.pathreplacing,backgrounds,shadows,babel}
\pgfplotsset{compat=1.18}
\usepgfplotslibrary{fillbetween}
\usepackage[np]{numprint}
\usepackage{listings,makecell,pas-tableur,twoopt}
\let\barre\undefined
\usepackage[pstricks]{tablvar}
\definecolor{darkgreen}{RGB}{0,100,0}
\providecommand{\red}{\color{red}}
\providecommand{\blue}{\color{blue}}
\providecommand{\green}{\color{green}}
\providecommand{\black}{\color{black}}
\providecommand{\white}{\color{white}}
\providecommand{\gray}{\color{gray}}
\providecommand{\cyan}{\color{cyan}}
\providecommand{\magenta}{\color{magenta}}
\providecommand{\yellow}{\color{yellow}}
\providecommand{\violet}{\color{violet}}
\providecommand{\orange}{\color{orange}}
\providecommand{\euro}{\texteuro}
\providecommand{\EUR}{\texteuro}
\providecommand{\ieme}{\textsuperscript{e}}
\providecommand{\degres}{\ensuremath{^\circ}}
\providecommand{\degree}{\ensuremath{^\circ}}
\providecommand{\vv}[1]{\overrightarrow{#1}}
\providecommand{\Coord}[1]{#1}
\providecommand{\decosix}{}
\providecommand{\starredbullet}{\ensuremath{\bullet}}
\pagestyle{empty}
\setlength{\parindent}{0pt}
\usepackage{adjustbox}\DeclareUnicodeCharacter{2236}{:}\DeclareUnicodeCharacter{2190}{\ensuremath{\leftarrow}}
\providecommand{\R}{\mathbb{R}}
\providecommand{\N}{\mathbb{N}}
\providecommand{\Z}{\mathbb{Z}}
\providecommand{\Q}{\mathbb{Q}}
\providecommand{\D}{\mathbb{D}}
\providecommand{\C}{\mathbb{C}}
\providecommand{\e}{{\mathrm{e}}}
\providecommand{\vect}[1]{\overrightarrow{#1}}
\providecommand{\vectt}[1]{\overrightarrow{\mathrm{#1}}}
\providecommand{\vvt}[1]{\overrightarrow{\mathrm{#1}}}
\providecommand{\np}[2][]{{#2~#1}}
\providecommand{\vv}[1]{{\overrightarrow{#1}}}
\renewcommand{\d}{{\mathrm{d}}}
\providecommand{\degres}{{^\circ}}
\providecommand{\degre}{{^\circ}}
\providecommand{\blue}{\color{blue}}
\providecommand{\red}{\color{red}}
\providecommand{\enskip}{\quad}
\providecommand{\llbracket}{\lbrack\!\lbrack}
\providecommand{\rrbracket}{\rbrack\!\rbrack}
\providecommand{\euro}{{\text{€}}}
\providecommand{\no}{{\text{n°}}}
\providecommand{\interval}[2]{{[#1;#2]}}
\providecommand{\intervFO}[2]{{[#1;#2[}}
\providecommand{\dotfill}{\ldots}

\providecommand{\ds}{\displaystyle}
\providecommand{\textding}[1]{#1}
\providecommand{\Oij}{\(\mathrm{O};\vec{\imath},\vec{\jmath}\)}
\providecommand{\Oijk}{\(\mathrm{O};\vec{\imath},\vec{\jmath},\vec{k}\)}
\providecommand{\Ouv}{\(\mathrm{O};\vec{u},\vec{v}\)}
\providecommand{\intervFF}[2]{\ensuremath{[#1;#2]}}
\renewcommand{\euro}{\ifmmode\text{\texteuro}\else\texteuro\fi}
\DeclareUnicodeCharacter{00B0}{\ifmmode^{\circ}\else\textdegree\fi}
\DeclareUnicodeCharacter{20AC}{\ifmmode\text{\texteuro}\else\texteuro\fi}
% Énoncés : Éducation nationale. Transcription et corrigés : APMEP, Denis Vergès et contributeurs. Découpage et classement : CoopMaths, Sébastien Lozano, Éric Elter et contributeurs. Adaptation de la présentation : Maths Explorer.
% CC BY-SA : https://coopmaths.fr/www/about/
% https://coopmaths.fr/alea/static/bac/2025/tex/bac_2025_05_sujet2_ameriquenord_3.tex
\begin{document}
\noindent\textit{APMEP / CoopMaths. CC BY-SA. Présentation adaptée par Maths Explorer.}\par\medskip
Pour chacune des affirmations suivantes, indiquer si elle est vraie ou fausse. Chaque réponse doit être justifiée.

Une réponse non justifiée ne rapporte aucun point.

\medskip

\textbf{PARTIE A}

\medskip

\begin{minipage}{0.55\linewidth}
ABCDEFGH est un cube d'arête de longueur 1.

Les points I, J, K, L et M sont les milieux respectifs des arêtes [AB], [BF], [AE], [CD] et [DH].

\medskip

\end{minipage}\hfill
\begin{minipage}{0.40\linewidth}
\psset{unit=1cm}
\begin{pspicture}(5.5,5.5)
\pspolygon(1.5,0.6)(5,0.9)(5,4.4)(1.5,4.1)
\psline(5,4.4)(4,5.4)(0.5,5.1)(1.5,4.1)
\psline(0.5,5.1)(0.5,1.6)(1.5,0.6)
\psline[linestyle=dashed](0.5,1.6)(4,1.9)(4,5.4)
\psline[linestyle=dashed](5,0.9)(4,1.9)
\uput[dl](1.5,0.6){A} \uput[dr](5,0.9){B} \uput[r](4,2){C} \uput[l](0.5,1.6){D}
\uput[l](1.5,4.1){E} \uput[r](5,4.4){F} \uput[ul](4,5.4){G} \uput[ul](0.5,5.1){H}
\uput[dr](3.25,0.75){I} \uput[r](5,2.65){J} \uput[l](1.5,2.35){K}\uput[u](2.25,1.75){L}
\uput[l](0.5,3.35){M}
\psdots(1.5,0.6)(5,0.9)(4,1.9)(0.5,1.6)(1.5,4.1)(5,4.4)(4,5.4)(0.5,5.1)(3.25,0.75)(5,2.65)(1.5,2.35)(2.25,1.75)(0.5,3.35)
\end{pspicture}
\end{minipage}

\smallskip

\textbf{Affirmation 1 :} \og $\vect{\text{JH}} = 2\vect{\text{BI}} + \vect{\text{DM}} - \vect{\text{CB}} $ \fg

\textbf{Affirmation 2 :} \og Le triplet de vecteurs $\left(\vect{\text{AB}},~\vect{\text{AH}},~\vect{\text{AG}}\right)$ est une base de l'espace. \fg

\textbf{Affirmation 3 :} \og $\vect{\text{IB}} \cdot \vect{\text{LM}} = - \dfrac 14$. \fg

\bigskip

\textbf{PARTIE B}

\medskip

Dans l'espace muni d'un repère orthonormé, on considère :

\begin{itemize}[label={$\bullet~$}]
\item le plan $\mathcal{P}$ d'équation cartésienne $2x - y + 3z + 6 = 0$
\item les points A$(2~;~0~;~-1)$ et B$(5~;~-3~;~7)$
\end{itemize}

\textbf{Affirmation 4 :} \og Le plan $\mathcal{P}$ et la droite (AB) sont parallèles.\fg

\textbf{Affirmation 5 :} \og Le plan $\mathcal{P}'$ parallèle à $\mathcal{P}$ passant par B a pour équation cartésienne $-2x + y - 3z + 34 = 0$ \fg

\textbf{Affirmation 6 :} \og La distance du point A au plan $\mathcal{P}$ est égale à $\dfrac{\sqrt{14}}{2}$.\fg

\smallskip

On note $(d)$ la droite de représentation paramétrique

\[\left\{\begin{array}{l c l}
x &=&-12 + 2k\\
y &=&\phantom{-}6\\
z &=&\phantom{-}3 - \phantom{-}5k
\end{array}\right., \text{où}\, \:k \in \R\]

\textbf{Affirmation 7 :} \og Les droites (AB) et $(d)$ ne sont pas coplanaires. \fg

\bigskip

\end{document}
